\subsection{Inference and Deployment Considerations}
\label{sec:inference}

At deployment time the trained model, the fitted scaler, and the label encoder are serialised to disk and loaded by the inference worker. Each incoming flow record is passed through the same \textsc{EngineerFeatures} and \textsc{Scale} steps that were applied at training time, ensuring that there is no training/serving skew. The trained classifier then emits a class probability vector; the predicted label is the class with the highest probability.

As shown in Section~\ref{sec:exp-latency}, all tree-based models achieve $> 30{,}000$ flows per second on a single CPU core, comfortably satisfying the $10^4$ flows/s throughput target. Micro-batching (collecting 128 flows before calling \texttt{predict\_proba} once) raises throughput by $6{-}9\times$ over single-sample calls with negligible additional latency, making batch inference the recommended default for production deployments. The serialised LightGBM model occupies only $\approx 5\,\text{MB}$, enabling co-deployment with the flow extractor on the same gateway host well within the 512 MB resource budget.
